\section{Conclusion}
\label{sec:conclusion}

The central question this paper set out to answer was whether a
P\&ID digitizer could be trained without access to any real
annotated plant drawings.
The answer, it turns out, depends almost entirely on how the synthetic
training data is made.
Random symbol placement --- the standard approach in prior work ---
produces diagrams whose graph topology bears little resemblance to real
process plants, and models trained on such data fail to transfer.
Seeding generation from real P\&IDs preserves the structural patterns
that matter, and the resulting \ours{} corpus closes most of the
gap to supervised training: 63.8\% edge mAP with zero real training
data, versus 69.4\% with the oracle and 45.9\% with the modular baseline.

The scaling analysis points to a clear next step.
Performance saturates around 400 images from 12 seeds, which
suggests that generating more images from the same seeds will not
help much.
Broadening the seed pool --- collecting even two or three dozen real
P\&IDs from different facilities --- is likely a more productive
use of effort than running the generator longer.
Together with the E3 result showing that a small amount of real data
significantly boosts performance, this suggests a practical two-stage
approach: generate a large synthetic corpus from whatever real P\&IDs
are available, then fine-tune briefly on a curated sample of real images.

We will release \ours{}, the generation pipeline, and trained models
upon acceptance in the hope that this provides a useful foundation
for future work on P\&ID digitization and, more broadly, on
synthetic-to-real transfer for structured diagram understanding.
